\documentclass[11pt]{article}
\usepackage[a4paper,margin=25mm]{geometry}
\usepackage{amsmath,amssymb,amsthm,array,booktabs}
\usepackage[T1]{fontenc}
\usepackage{lmodern}
\usepackage[hidelinks]{hyperref}
\newtheorem{theorem}{Theorem}
\newtheorem{corollary}[theorem]{Corollary}
\newcommand{\den}{\operatorname{den}}
\newcommand{\num}{\operatorname{num}}
\title{Catalan formula corrections and short proofs\\of Groux's and Staunton's conjectures}
\author{Haobo Ma\\\small ChronoAI Pte Ltd, Singapore
\and Wenlin Zhang\\\small National University of Singapore\\\small ChronoAI Pte Ltd, Singapore}
\date{November 13, 2026}
\begin{document}
\maketitle
\begin{abstract}
Applying Pieter Mostert's finite-sum identity (OEIS A119951,
October 12, 2012) and the minimal-exponent theorem of Ca\c{c}\~ao,
Falc\~ao and Malonek (2017, Theorem 6), we give short proofs of
Groux Roland's denominator conjecture and Alexander Staunton's
linked exponent conjecture in the FORMULA section of OEIS A120069
(December 11, 2010 and February 28, 2026, respectively).
For every positive index, the reduced denominator of the partial sum
of scaled Catalan numbers equals that of its last summand; we also
give its binary-digit form. We correct an index in the companion
numerator entry and a minimum expression in the exponent entry.
\end{abstract}

\section{The denominator equality}
Write $C_n=\binom{2n}{n}/(n+1)$ for the Catalan number, and put
\[
 q_n=\frac{C_n}{4^{n-1}},\qquad
 S_n=\sum_{k=1}^n q_k,\qquad a_n=\den(S_n)\quad(n\geq1).
\]
Here $\den$ is the positive denominator in lowest terms, with
$\den(0)=1$. The sequences of numerators and denominators of $S_n$
are A119951 and A120069 respectively \cite{numerators,groux}.
The point of Groux's conjecture is that adding all earlier summands
does not change the reduced denominator of the final one.

\begin{theorem}\label{thm:main}
For every integer $n\geq1$, $a_n=\den(q_n)$.
\end{theorem}
\begin{proof}
The factorial definition gives
\[
 (n+2)C_{n+1}=2(2n+1)C_n.
\]
Consequently $(2n+3)q_{n+1}+q_{n+1}=(2n+1)q_n$.
Starting with $S_1=q_1=1$, induction gives
\begin{equation}\label{eq:mostert}
 S_n=4-(2n+1)q_n
     =4-\frac{\binom{2n+2}{n+1}}{2^{2n-1}}.
\end{equation}
This is Mostert's identity, recorded in A119951 on October 12, 2012
\cite{numerators}; the induction supplies its proof here.
Write $q_n=u/d$ in lowest terms. Since $d\mid4^{n-1}$, $d$ is a power
of two, and $\gcd(2n+1,d)=1$. Thus
\[
 S_n=\frac{4d-(2n+1)u}{d},\qquad
 \gcd\bigl(4d-(2n+1)u,d\bigr)=1.
\]
Its reduced denominator is therefore exactly $d$, including $d=1$.
\end{proof}

\section{The linked exponent conjecture}
Let $b_n=\mathrm{A283208}(n)$ for $n\geq0$. The definition in that
entry \cite{vietoris} is equivalently
\begin{equation}\label{eq:bdef}
 2^{b_n}=\den\left(\frac{(2n+1)C_n}{4^n}\right).
\end{equation}
Indeed, its Vietoris number $c(2n)$ is $\binom{2n}{n}/4^n$,
and $b_n$ is the least nonnegative exponent making
$2^{b_n}(2n+1)c(2n)/(n+1)$ integral. This matches the
minimal-exponent convention, not a shifted sequence index.
The rational in \eqref{eq:bdef} is $T_0^n$ in
Ca\c{c}\~ao, Falc\~ao and Malonek \cite[Eq.~(5.6)]{cacao2017}.
Their Theorem~6 proves its reduced denominator exponent, namely
\eqref{eq:factorial} below. Together with Mostert's identity this
already implies both A120069 conjectures. The following proof
rederives the exponent and gives its binary-digit form.

\begin{corollary}\label{cor:staunton}
Staunton's conjecture is true: for every $n\geq1$,
\[
 a_n=2^{b_n-2}.
\]
More explicitly, if $s_2(m)$ denotes the number of ones in the binary
expansion of $m$, then
\begin{equation}\label{eq:digits}
 a_n=2^{\,2n-1-s_2(n+1)},\qquad
 b_n=2n+1-s_2(n+1).
\end{equation}
The displayed formula for $b_n$ also holds at $n=0$.
\end{corollary}
\begin{proof}
Let $v_2(m)$ denote the exponent of two in a positive integer $m$.
Counting multiples of each power of two in $m!$ gives
\[
 v_2(m!)=\sum_{j\geq1}\left\lfloor\frac{m}{2^j}\right\rfloor
        =m-s_2(m).
\]
The second equality follows by writing $m=\sum_i\epsilon_i2^i$:
the sum of floors is $\sum_i\epsilon_i(2^i-1)$.
Applying this to $C_n$ yields
\[
 v_2(C_n)=s_2(n)-v_2(n+1)=s_2(n+1)-1.
\]
For the last equality, adding one to $n$ deletes precisely
$v_2(n+1)$ trailing ones and inserts one new one.
For $n=1$, $C_1=1$; for $n\geq2$, the Catalan recurrence and
$C_2=2<4$ give $0<C_n<4^{n-1}$, since $C_{n+1}/C_n<4$.
Hence $v_2(C_n)\leq2n-2$ for all $n\geq1$.
Reduction of $C_n/4^{n-1}$ and of the rational in
\eqref{eq:bdef} therefore gives, respectively,
\[
 a_n=2^{2n-2-v_2(C_n)},\qquad b_n=2n-v_2(C_n).
\]
Theorem~\ref{thm:main} and the valuation identity prove all claims.
At $n=0$, \eqref{eq:bdef} gives $b_0=0$.
\end{proof}

\paragraph{Other formulas explained by the same valuation.}
For $n\geq0$ we also obtain the existing A283208 formula,
recorded by Ca\c{c}\~ao, Malonek, Falc\~ao and Tomaz
\cite[Section 2, pp.~4 and 7]{cacao2018},
\begin{equation}\label{eq:factorial}
 b_n=n+v_2((n+1)!)
     =n+\sum_{j\geq1}\left\lfloor\frac{n+1}{2^j}\right\rfloor,
\end{equation}
and hence $b_{n+1}-b_n=1+v_2(n+2)$.
This proves both alternating identities recorded by Rog\'erio
Ser\^odio (February 19, 2019) \cite{vietoris}:
\[
 -b_{4n-1}-b_{4n}+b_{4n+1}+b_{4n+2}=6\quad(n\geq1),
\]
\[
 b_{2n+1}-b_{2n+2}-b_{2n+3}+b_{2n+4}=0\quad(n\geq0).
\]
In the first, group the differences two indices apart; each is three.
In the second, the two first differences being subtracted both equal one.
Similarly, with $m=\lceil(n+2)/4\rceil$,
\[
 b_{n+4}-b_n=7+v_2(m)\quad(n\geq0),
\]
because the four consecutive integers $n+2,\ldots,n+5$ contain
one odd pair, one integer of valuation one and the integer $4m$.
This is Ser\'odio's recurrence $6+\mathrm{A001511}(m)$.
The alternative minimum expression in A283208, read literally as
$\min\{k:\lfloor((n+5)/4)\bmod2^k\rfloor=1\}$, needs a correction:
at $n=3$ its residue is zero for $k=1$ and two for every $k\geq2$,
so that minimum does not exist. A precise replacement for the
increment beyond six is
$\min\{k\geq1:m\bmod2^k>0\}=1+v_2(m)$.
The empirical bitwise formula attributed to Gary Detlefs (2019)
and marked true by Andrei Zabolotskii (2021) is also
\eqref{eq:factorial}: the least set bit of $M=(n+1)!$ equals
$M-(M\mathbin{\mathrm{AND}}(M-1))=2^{v_2(M)}$.
These are rederivations of formulas already present in A283208.
For completeness its displayed generating function follows as well.
Writing $B(x)=\sum_{n\geq0}b_nx^n$ and using
$v_2(m)=\sum_{j\geq1}\mathbf{1}_{2^j\mid m}$ in the four-step
recurrence gives the formal power-series identity
\[
 (1-x^4)B(x)=\frac{2x+x^2+3x^3+x^4}{1-x}
 +\sum_{j\geq0}\frac{\sum_{i=0}^3x^{8\cdot2^j-1+i}}{1-x^{8\cdot2^j}}.
\]
Indeed, $b_0,b_1,b_2,b_3=0,2,3,6$ give the initial terms;
the constant seven in the recurrence gives every coefficient from
$x^4$ onward in the first fraction. For $r\geq1$, the condition
$2^r\mid\lceil(n+2)/4\rceil$ contributes blocks at degrees
$4\cdot2^r\ell-1,\ldots,4\cdot2^r\ell+2$, $\ell\geq1$,
which give the displayed sum after putting $r=j+1$.

The numerator $U_n=\num(S_n)$ has the companion expression
\[
 U_n=4a_n-(2n+1)\frac{C_n}{2^{v_2(C_n)}}.
\]
It follows directly from \eqref{eq:mostert}. Substitution of
\eqref{eq:digits} gives the equivalent factorial expression
attributed to Detlefs (November 6, 2020) in A119951:
\[
 U_n=\left(2^n-\frac{\binom{2n+2}{n+1}}{2^{n+1}}\right)
       \gcd((n+1)!,2^{n+1}).
\]
Here the greatest common divisor is
$2^{n+1-s_2(n+1)}$, since $v_2((n+1)!)\leq n$.

\paragraph{An indexing correction in A119951.}
The FORMULA line attributed to Johannes W. Meijer (May 24, 2009)
uses the numerator of a central-binomial convolution for $a(n-1)$.
With that entry's offset one, the correct index is $a(n)$.
To see this, put $H_j=\binom{2j}{j}$. The formal binomial series
$\sum_{j\geq0}H_jx^j=(1-4x)^{-1/2}$ gives
$\sum_{j=0}^{n+1}H_jH_{n+1-j}=4^{n+1}$. Removing the two endpoints
and using \eqref{eq:mostert} proves, for $n\geq1$,
\[
 \frac{1}{4^n}\sum_{i=0}^{n-1}H_{i+1}H_{n-i}
 =4-\frac{2H_{n+1}}{4^n}=S_n.
\]
Thus its numerator is $U_n$, not $U_{n-1}$.
For example, $n=2$ gives $24/16=3/2$, whose numerator is three,
whereas $U_1=1$. The convolution programs in the same entry
already produce the correctly indexed sequence.

\section{Coverage and verification}
\begin{center}\small
\begin{tabular}{@{}p{.46\linewidth}p{.46\linewidth}@{}}
\toprule Source statement & Coverage \\ \midrule
A120069, Groux (2010) & Proved from Mostert's identity (2012), Theorem~\ref{thm:main}. \\
A119951, Mostert (2012) & Reproved; Mostert (2012), \eqref{eq:mostert}. \\
A120069, Staunton (2026) & Proved as a consequence of Mostert (2012) and Ca\c{c}\~ao--Falc\~ao--Malonek (2017, Theorem~6), Corollary~\ref{cor:staunton}. \\
A283208, floor-sum exponent & Reproved; prior proof Ca\c{c}\~ao--Falc\~ao--Malonek (2017, Theorem~6). \\
A283208, Ser\'odio identities, recurrence and generating function; Detlefs bitwise formula & Rederived above; existing formulas credited. \\
A283208, minimum-expression version of the recurrence & Refuted-and-corrected by the positive-residue minimum above. \\
A119951, Detlefs numerator formula (2020) & Rederived above. \\
A119951, Meijer convolution formula (2009) & Refuted-and-corrected: replace $a(n-1)$ by $a(n)$. \\
A119951, Eta-triangle connections & Not covered. \\
\bottomrule
\end{tabular}
\end{center}
The pinned Lean file \cite{lean} contains the declaration
\begin{center}\small
\texttt{D5.S3.ArithSums.GrouxA120069Denominator.result}
\end{center}
Its statement is Theorem~\ref{thm:main} for natural $n$ with $0<n$.
Its private helpers \texttt{balance} and \texttt{sum\_formula}
establish the Catalan recurrence and \eqref{eq:mostert}.
The valuation formulas, Staunton corollary, numerator formula and
related recurrences and the convolution correction have the written proofs given above.
The accompanying \texttt{verify.py} checks the exact rational identities
for $1\leq n\leq1000$ and the displayed A283208 difference formulas
over their stated finite test ranges; these checks supplement the proofs.

\begin{thebibliography}{9}
\bibitem{groux} OEIS Foundation Inc., \emph{A120069},
revision 30, March 10, 2026. FORMULA: Groux Roland, December 11, 2010;
Alexander Staunton, February 28, 2026.
\url{https://oeis.org/A120069}.
\bibitem{numerators} OEIS Foundation Inc., \emph{A119951},
revision 32, February 16, 2025. COMMENT: Pieter Mostert,
October 12, 2012; FORMULA: Gary Detlefs, November 6, 2020.
\url{https://oeis.org/A119951}.
\bibitem{vietoris} I. Ca\c{c}\~ao, M. I. Falc\~ao and H. Malonek,
\emph{A283208}, OEIS, revision 103, August 5, 2026.
Additional formulas by R. Ser\'odio, G. Detlefs and A. Zabolotskii
are individually dated in the entry.
\url{https://oeis.org/A283208}.
\bibitem{cacao2017} I. Ca\c{c}\~ao, M. I. Falc\~ao and H. R. Malonek,
\emph{Hypercomplex Polynomials, Vietoris' Rational Numbers and a
Related Integer Numbers Sequence}, Complex Analysis and Operator Theory
\textbf{11} (2017), 1059--1076, Theorem~6 and Eq.~(5.6).
\url{https://doi.org/10.1007/s11785-017-0649-5}.
\href{https://hdl.handle.net/1822/45794}{Author manuscript}.
\bibitem{cacao2018} I. Ca\c{c}\~ao, H. R. Malonek, M. I. Falc\~ao
and G. Tomaz, \emph{Combinatorial Identities Associated with a
Multidimensional Polynomial Sequence}, Journal of Integer Sequences
\textbf{21} (2018), Article 18.7.4.
\url{https://cs.uwaterloo.ca/journals/JIS/VOL21/Falcao/falcao2.pdf}.
\bibitem{lean} \emph{GrouxA120069Denominator.lean},
commit \texttt{3378450a2176909e857354ab064dde6ec7679b71}.
\href{https://github.com/the-omega-institute/trureturing/blob/3378450a2176909e857354ab064dde6ec7679b71/D5/S3/ArithSums/GrouxA120069Denominator.lean}{Pinned Lean source}.
\end{thebibliography}
\bigskip\noindent\small
Copyright \textcopyright\ 2026 Haobo Ma and Wenlin Zhang.
Licensed under Creative Commons Attribution 4.0 International (CC BY 4.0).
\end{document}
